% TOP_PROBLEM: {"id": 2302059, "title": "Research Problems in Function Theory — Problem 2.59", "category_id": 8, "category": {"id": 8, "name": "analysis", "display_name": "Analysis", "description": "Limits, continuity, calculus, and function theory.", "slug": "analysis", "order_index": 8, "created_at": "2026-07-31T15:26:25.671Z"}, "status": "partially_solved", "problem_number": "AMR-022-2059", "set_id": 13}
% FIRST_POSTED: 2026-10-02
% ENTRY_KIND: statement_only
% UnsolvedMath ID 2302059; source catalogue code AMR-022-2059
% !TeX program = lualatex
% Definition and sources only; this file is not an LLM solution attempt.
\documentclass[11pt,a4paper]{article}
\usepackage[margin=25mm]{geometry}
\usepackage{fontspec}
\setmainfont{lmroman10-regular.otf}[BoldFont=lmroman10-bold.otf,
 ItalicFont=lmroman10-italic.otf,BoldItalicFont=lmroman10-bolditalic.otf]
\usepackage{amsmath,amssymb,mathtools}
\usepackage{unicode-math}
\setmathfont{latinmodern-math.otf}
\AtBeginDocument{\let\setminus\smallsetminus}
\usepackage{xurl,hyperref}
\hypersetup{colorlinks=true,urlcolor=blue}
\setcounter{secnumdepth}{0}
\setlength{\parindent}{0pt}
\setlength{\parskip}{5pt}
\setlength{\emergencystretch}{3em}
\begin{document}
\section*{Research Problems in Function Theory — Problem 2.59}\label{2302059}
{\small UnsolvedMath ID 2302059 · AMR-022-2059 · Analysis · AMR Open Problem Lists}
\subsection{Definitions and mathematical statement}
Let $\rho\ge0$ and $K>0$. Consider a complex power series with coefficients $(a_k)_{k\ge0}$, and write its sections as
$S_n(z)=\sum_{k=0}^na_kz^k$ for $n\ge1$.
Assume every section is nonzero throughout the same open region
\[
W_{\rho,K}=\{x+iy:x>0,\ |y|<Kx^{1-\rho/2}\}.
\]
In particular no section is identically zero. The width conjecture asks whether this hypothesis forces the series to have infinite radius of convergence, with its sum $f$ satisfying
\[
\operatorname{ord}(f)=\limsup_{r\to\infty}
\frac{\log\log M(r,f)}{\log r}\le\rho,
\qquad M(r,f)=\max_{|z|=r}|f(z)|.
\]
Polynomials are taken to have order zero. If no convergence is initially assumed, infinite radius is itself part of the conclusion; if the series already converges near zero, the same assertion still includes continuation by that Taylor series to the whole plane.
The region has angular width when $\rho=0$ and narrows at the rate dictated by the exponent otherwise. Both $K$ and $\rho$ are independent of $n$. A different zero-free region for each section, or zero-freeness only for a subsequence of sections, is not the stated hypothesis.
\subsection{Short English statement}
Does a common zero-free region of the prescribed width for every Taylor section force an entire sum of order at most the associated exponent?
\subsection{Sources}
\begin{itemize}
\item[S1] ULAM.AI and the UnsolvedMath Contributors, supplied problems.json, record 2302059 (AMR-022-2059), AMR Open Problem Lists. Catalogue identity and source transcription; CC BY 4.0.
\url{https://huggingface.co/datasets/ulamai/UnsolvedMath}.
\item[S2] Hayman - Research Problems in Function Theory (2018), Problem 2.59. Original source indexed by AMR-022-2059.
\url{https://arxiv.org/abs/1809.07200}.
\end{itemize}
\end{document}
